\chapter{Implementation}
\label{ch:implementation}

\section{Frontend Implementation}
The frontend is built with React and bundled with Vite, using plain
JavaScript (no TypeScript) and component-scoped CSS. No external UI
component framework is used for the game board itself; the board is
rendered as a CSS grid of cell components, each responding to left-click
(reveal), right-click (flag), and, on an already-revealed numbered cell,
click-again (chord).

\section{Backend Implementation}
The backend is written in procedural PHP, without a framework, as a
collection of standalone endpoint scripts plus shared includes for the
database connection and authentication helpers. Unlike a typical
\texttt{mysqli}-based implementation, database access uses PHP's PDO
extension with prepared statements throughout, so that every query
involving user-supplied data is parameterized.

\section{API Endpoints}
Table~\ref{tab:api-endpoints} summarises the backend endpoints exposed by
the application.

\begin{table}[h]
\centering
\caption{Summary of backend API endpoints.}
\label{tab:api-endpoints}
\begin{tabular}{@{}p{3.4cm}p{1.6cm}p{1.8cm}p{6.5cm}@{}}
\toprule
\textbf{Endpoint} & \textbf{Method} & \textbf{Auth} & \textbf{Purpose} \\
\midrule
register.php & POST & No & Create a new user account \\
login.php & POST & No & Authenticate and start a session \\
logout.php & POST & Yes & Destroy the current session \\
save-score.php & POST & Yes & Recompute and persist a completed game's result \\
leaderboard.php & GET & Yes & Return the top scores, optionally filtered by difficulty \\
history.php & GET & Yes & Return the logged-in user's past games \\
stats.php & GET & Yes & Return the logged-in user's aggregate statistics \\
\bottomrule
\end{tabular}
\end{table}

\section{Representative Frontend Code: API Wrapper}
\begin{lstlisting}[language=JavaScript, caption={Simplified fetch wrapper (\texttt{services/api.js}).}, label={lst:api-js}]
const API_BASE = import.meta.env.VITE_API_URL;

export async function apiRequest(endpoint, options = {}) {
  const response = await fetch(`${API_BASE}${endpoint}`, {
    credentials: "include",
    headers: { "Content-Type": "application/json" },
    ...options,
  });
  const data = await response.json();
  if (!response.ok || !data.success) {
    throw new Error(data.message || "Request failed");
  }
  return data;
}
\end{lstlisting}

\section{Representative Backend Code: Score Submission Endpoint}
\texttt{save-score.php} is the most security-relevant endpoint in the
system, since it is the only point at which client-reported gameplay
metrics are converted into a persisted, leaderboard-visible score. Its
validation logic is described fully in Section~\ref{sec:score-validation}
of Chapter~\ref{ch:security}; a simplified outline is shown below.

\begin{lstlisting}[language=PHP, caption={Simplified score submission endpoint (\texttt{save-score.php}).}, label={lst:save-score}]
<?php
require_once "includes/db.php";
require_once "includes/auth.php";
require_login();

$input = json_decode(file_get_contents("php://input"), true);
$difficulty = $input["difficulty"];
$timeTaken = (int) $input["time_taken"];
$cellsRevealed = (int) $input["cells_revealed"];
$hintsUsed = (int) $input["hints_used"];

// Bound-check against the known limits for this difficulty
validate_metrics($difficulty, $timeTaken, $cellsRevealed, $hintsUsed);

// Reject if an identical result was submitted in the last 10 seconds
reject_if_duplicate($userId, $difficulty, $timeTaken, $cellsRevealed, $hintsUsed);

// Recompute the score server-side; never trust a client-supplied score
$score = max(0, $cellsRevealed * 10 - $hintsUsed * 50);

$stmt = $pdo->prepare(
  "INSERT INTO game_results (user_id, difficulty, score, time_taken, cells_revealed, hints_used)
   VALUES (?, ?, ?, ?, ?, ?)"
);
$stmt->execute([$userId, $difficulty, $score, $timeTaken, $cellsRevealed, $hintsUsed]);
\end{lstlisting}

\section{Configuration Management}
Difficulty parameters are defined once on the frontend in the
\texttt{DIFFICULTIES} configuration object (grid size and mine count) and
independently on the backend as the bounds used by
\texttt{validate\_metrics()} (maximum safe cells per difficulty). This
duplication -- rather than a single shared source of truth -- is noted as a
maintainability limitation in Chapter~\ref{ch:limitations}.
